% GENERATED FILE. Edit canonical lessons, then run npm run generate:books.
% canonical-source-hash: fnv1a64:af61b04e03153ec1
% canonical-lessons: ES-C302-calle, ES-C302-esquina, ES-C302-huella, ES-C302-puerto, ES-C302-faro

\chapter{Out and Away --- Street, Corner, Footprint, Harbour, Lighthouse}
\label{ch:calle-y-faro}
\hlchaptermodality{302}

\begin{chapteropening}
\textbf{By the end of this chapter:} \emph{I can name the street, the corner, the footprint, the harbour and the lighthouse in Spanish.}

Complete ES-C302-faro: name all five of a leaving in order, from the street to the light behind you.
\end{chapteropening}

\section[la calle]{la calle — a cattle track that reached the town}

\label{lesson:ES-C302-calle}



\begin{quote}
Before a new run: which English word is \emph{cista}, worn down? (\emph{Chest}.) And what plant is \emph{papel} named after? (The papyrus reed.)
\end{quote}

\subsection*{What to know first}
\begin{itemize}
\raggedright
  \item el sendero — the way through the trees.
  \item el ganado — what wore this one flat in the first place.
\end{itemize}

\begin{sounds}
\begin{itemize}
\raggedright
  \item \emph{CA-lle}, stress on the first part. The \emph{c} is hard and the \emph{ll} is a \emph{y}. $\to$ reference
\end{itemize}
\end{sounds}

\begin{cousinweb}
\textbf{la calle} — \textquotedblleft{}the street.\textquotedblright{}

Latin \textbf{callis} was not urban at all. It was the narrow hard track that animals beat through rough country when they were driven from one grazing ground to another — a hoof-path, made by use rather than by anybody's decision.

Then towns grew, and the word walked in with the herds. Spanish now uses \emph{calle} for the most built thing there is: a paved street with houses down both sides. The stones are new; the word underneath them was made by \emph{el ganado}.

You have three words for a way through now, and they do not overlap. \emph{El sendero} is the path through the trees, narrow and unpaved. \emph{El camino} is the road anybody would take. \emph{La calle} has walls on both sides.

Ask \emph{¿en qué calle vives?} and you are asking which street somebody lives on.
\end{cousinweb}

\begin{grammarlens}[title={What you've built}]
The first step of a leaving, and it starts at the door.

Everything that follows in this run is further out than the last: a street, then its corner, then what you leave in the dirt beyond it.
\end{grammarlens}

\subsection*{Your turn}
Say these aloud:

\begin{itemize}
\raggedright
  \item \emph{la calle}
  \item \emph{el sendero}, then \emph{la calle}
  \item \emph{¿en qué calle vives?}
\end{itemize}

\begin{tcolorbox}[breakable,colback=teal!4,colframe=teal!35!black,title={Before you move on}]
What was a \emph{callis}? (A track beaten by driven animals.) What is \emph{la calle} now? (A street.) And which of your three ways through has walls on both sides? (\emph{La calle}.)
\end{tcolorbox}
\section[la esquina]{la esquina — named after a shinbone}

\label{lesson:ES-C302-esquina}



\begin{quote}
Two back: what does \emph{papel} also mean, besides paper? (The part an actor plays.) And what was a \emph{callis}? (A track beaten by animals.)
\end{quote}

\subsection*{What to know first}
\begin{itemize}
\raggedright
  \item la calle — what this is the end of.
  \item la escoba — another word wearing the same \emph{e-} step.
\end{itemize}

\begin{sounds}
\begin{itemize}
\raggedright
  \item \emph{es-QUI-na}, stress in the middle. The \emph{qui} is a hard \emph{k} with the \emph{u} silent. $\to$ reference
\end{itemize}
\end{sounds}

\begin{cousinweb}
\textbf{la esquina} — \textquotedblleft{}the corner,\textquotedblright{} the outside one, where two walls meet.

This one is not Latin. It came from the Germanic speakers, whose \textbf{skina} meant a shinbone, or a thin splinter of bone or wood — a long hard edge. A corner is exactly that shape, and the language borrowed the bone to name it.

English kept the bone and did not lend it to the corner: your \textbf{shin} is \emph{skina}, the same word, and anybody who has met a table corner in the dark can confirm that the two ideas belong together.

Notice the front of the Spanish word. \emph{Skina} could not stand as it was, because Spanish does not begin a word with \emph{s} and another consonant — so it put an \emph{e} on the front, exactly as it did to make \emph{escoba} and \emph{estar}. The \emph{e-} is not part of the word. It is a step Spanish takes to get onto it.
\end{cousinweb}

\begin{grammarlens}[title={What you've built}]
Two of the leaving, and a step you can now recognise on sight.

Any Spanish word starting \emph{es-} with a consonant after it is worth checking: strip the \emph{e}, and a Latin or Germanic word is often standing underneath.
\end{grammarlens}

\subsection*{Your turn}
Say these aloud:

\begin{itemize}
\raggedright
  \item \emph{la esquina}
  \item \emph{la esquina de la calle}
  \item \emph{la escoba}, then \emph{la esquina} — the same step twice
\end{itemize}

\begin{tcolorbox}[breakable,colback=teal!4,colframe=teal!35!black,title={Before you move on}]
What did \emph{skina} mean? (A shinbone, or a splinter.) Which English word is it still? (\emph{Shin}.) And why does \emph{esquina} start with an \emph{e}? (Spanish will not begin a word with \emph{s} plus a consonant.)
\end{tcolorbox}
\section[la huella]{la huella — what the treading leaves}

\label{lesson:ES-C302-huella}



\begin{quote}
What was a \emph{callis}? (A track beaten by driven animals.) And what did \emph{skina} mean? (A shinbone.)
\end{quote}

\subsection*{What to know first}
\begin{itemize}
\raggedright
  \item el hilo — the last word this run gave you with a buried \emph{f}.
  \item el sendero — a thing made entirely out of these.
\end{itemize}

\begin{sounds}
\begin{itemize}
\raggedright
  \item \emph{UE-lla}, and the \emph{h} is silent, so the word opens straight onto the \emph{ue} slide. $\to$ reference
\end{itemize}
\end{sounds}

\begin{cousinweb}
\textbf{la huella} — \textquotedblleft{}the footprint,\textquotedblright{} \textquotedblleft{}the track,\textquotedblright{} the mark a thing leaves by pressing on something.

Third time for the decoder, and it still works. Silent \emph{h}, so put the \emph{f} back: behind \emph{huella} is \emph{hollar}, to tread, and behind that Latin \textbf{fullare}.

\emph{Fullare} was a specific job. Newly woven wool had to be cleaned and thickened, and the way to do it was to put it in a vat with water and earth and \textbf{tread on it}, for hours. The man who did it was a \emph{fullo}. English still carries him: to \textbf{full} cloth is a real word, a \textbf{fulling mill} did it with hammers once water power arrived, and the surname \textbf{Fuller} belonged to whoever did it by foot.

So a footprint in Spanish is named after a day's work in a vat of wet wool. And the word has gone abstract in both languages the same way: \emph{la huella} is also the mark somebody leaves behind them in general, exactly as English uses \textbf{footprint} for something that never had a foot.
\end{cousinweb}

\begin{grammarlens}[title={What you've built}]
Three of the leaving, and three uses of one decoder in two runs.

\emph{Horno}, \emph{hilo}, \emph{huella}. Each of them was an \emph{f} once, and each of them gave the letter up in the same century. A silent \emph{h} is now a piece of evidence, not a spelling to memorise.
\end{grammarlens}

\subsection*{Your turn}
Say these aloud:

\begin{itemize}
\raggedright
  \item \emph{la huella}
  \item \emph{el hilo}, then \emph{la huella} — two buried \emph{f} sounds
  \item \emph{la huella en el sendero}
\end{itemize}

\begin{tcolorbox}[breakable,colback=teal!4,colframe=teal!35!black,title={Before you move on}]
What letter is under the \emph{h} of \emph{huella}? (An \emph{f}.) What was \emph{fullare}? (Treading cloth in a vat.) And what does \emph{la huella} mean? (A footprint, or a mark left behind.)
\end{tcolorbox}
\section[el puerto]{el puerto — the other half of a pair}

\label{lesson:ES-C302-puerto}



\begin{quote}
Three questions. What was a \emph{callis}? (A beaten track.) Which English word is \emph{skina}? (\emph{Shin}.) And what was \emph{fullare}? (Treading cloth in a vat.)
\end{quote}

\subsection*{What to know first}
\begin{itemize}
\raggedright
  \item la puerta — the twin of this word, met long ago.
  \item la orilla — the edge this is built on.
\end{itemize}

\begin{sounds}
\begin{itemize}
\raggedright
  \item \emph{PUER-to}, stress on the first part. The \emph{ue} is one slide and the \emph{r} is a single tap. $\to$ reference
\end{itemize}
\end{sounds}

\begin{cousinweb}
\textbf{el puerto} — \textquotedblleft{}the harbour,\textquotedblright{} \textquotedblleft{}the port.\textquotedblright{}

When you took \emph{la puerta} apart, the Latin behind it was \textbf{porta}, a city gate. Latin had a second word beside it, built on the same idea of a way through: \textbf{portus}, the place a ship gets in. Both broke their stressed \emph{o} into \emph{ue}, and Spanish ended up with a matched pair — \emph{puerta} for the gate on land, \emph{puerto} for the gate off the sea.

English took \emph{portus} and never stopped using it. A \textbf{port} is the harbour. To \textbf{import} is to carry in through one; to \textbf{export}, out. \textbf{Portugal} is \emph{Portus Cale}, the harbour at Cale, which is where the whole country got its name.

And the best of them: \textbf{opportunity}. Latin \emph{ob portum} described a wind standing toward the harbour — the wind that lets you get in. An opportunity is not luck. It is the wind being in the right quarter, and somebody having a ship.
\end{cousinweb}

\begin{grammarlens}[title={What you've built}]
Four of the leaving, and a word you half owned already.

\emph{Puerta} and \emph{puerto} differ by one letter and by which side of the water you are standing on. Pairs like that are cheap to hold and easy to hear, and Spanish has more of them than any list would suggest.
\end{grammarlens}

\subsection*{Your turn}
Say these aloud:

\begin{itemize}
\raggedright
  \item \emph{el puerto}
  \item \emph{la puerta}, then \emph{el puerto} — the gate and the harbour
  \item all four — \emph{la calle}, \emph{la esquina}, \emph{la huella}, \emph{el puerto}
\end{itemize}

\begin{tcolorbox}[breakable,colback=teal!4,colframe=teal!35!black,title={Before you move on}]
What was \emph{portus}? (A harbour.) What is \emph{Portugal} named after? (The harbour at Cale.) And what does \emph{ob portum} describe? (A wind standing toward harbour — an opportunity.)
\end{tcolorbox}
\section[el faro]{el faro — an island that became a word}

\label{lesson:ES-C302-faro}



\begin{quote}
Before the last of the five: what was a \emph{callis}? (A beaten track.) What is under the \emph{h} of \emph{huella}? (An \emph{f}.) And what does \emph{ob portum} describe? (A wind standing toward harbour.)
\end{quote}

\subsection*{What to know first}
\begin{itemize}
\raggedright
  \item el puerto — what this stands at the mouth of.
  \item la isla — what the original one was built on.
\end{itemize}

\begin{sounds}
\begin{itemize}
\raggedright
  \item \emph{FA-ro}, stress on the first part. The \emph{r} is one tap, not a roll. $\to$ reference
\end{itemize}
\end{sounds}

\begin{cousinweb}
\textbf{el faro} — \textquotedblleft{}the lighthouse.\textquotedblright{}

There is no root to take apart, because this word was a place. \textbf{Pharos} was a small island off Alexandria, and on it the Ptolemies built a tower with a fire at the top that stood for some sixteen hundred years. Sailors called any such tower a \emph{pharos} after it, and the name swallowed the object: French \emph{phare}, Italian \emph{faro}, Spanish \emph{faro}.

The tower fell in an earthquake and the island is now joined to the mainland. Nothing of it is left except the word — which is also, in modern Spanish, what you call a car's headlights. The last of the great wonders is switched on every evening in every car park in Spain.
\end{cousinweb}

\begin{grammarlens}[title={What you've built}]
Five: \emph{la calle}, \emph{la esquina}, \emph{la huella}, \emph{el puerto}, \emph{el faro}.

Read in order they are a leaving that gets further out at every step — the street, its corner, the mark left in the dirt, the water's edge, and the light still visible when the town is not.

The five came from five different quarries: a cattle track, a shinbone, a day spent treading wool, a Roman harbour, and an island with a fire on it. That is why taking words apart keeps paying — you never know which one will turn out to be a building you have heard of.
\end{grammarlens}

\subsection*{Your turn}
Say these aloud:

\begin{itemize}
\raggedright
  \item \emph{el faro}
  \item \emph{el puerto}, then \emph{el faro}
  \item all five — \emph{la calle}, \emph{la esquina}, \emph{la huella}, \emph{el puerto}, \emph{el faro}
\end{itemize}

\begin{tcolorbox}[breakable,colback=teal!4,colframe=teal!35!black,title={Before you move on}]
What was Pharos? (An island off Alexandria with a great lighthouse on it.) How did it become the word? (Sailors called every such tower after it.) And what else is a \emph{faro} today? (A car's headlight.)
\end{tcolorbox}
